\begin{remarks}[5.9.12]
\label{IV.2.s5.9.12-fr}
\begin{enumerate}
\item[(i)] Sea $\boldsymbol{C}(X)$ la categoría de los $\mathcal O_X$-módulos y
sea $\boldsymbol{C}_Z(X)$ la subcategoría de $\boldsymbol{C}(X)$ formada por los
$\mathcal O_X$-módulos cuyo soporte está contenido en $Z$; esta subcategoría es
localizante en el sentido de Gabriel, y el funtor $\mathcal H^0_{X/Z}$ no es
sino el funtor de localización de Gabriel (cf. [27]; esto proporcionaría otra
demostración de (5.9.11)). Cuando $Z$ es cerrado, el funtor
$i^*:\boldsymbol{C}(X)\to\boldsymbol{C}(X\setmin Z)$ (donde $i:X\setmin Z\to X$ es la inyección
canónica) define una \emph{equivalencia} de categorías
$\boldsymbol{C}(X)/\boldsymbol{C}_Z(X)\approx\boldsymbol{C}(X\setmin Z)$.
\item[(ii)] De (5.9.11) resulta que la condición
$\mathcal H^0_{X/Z}(\mathcal F)=0$ es \emph{equivalente} a
$\operatorname{Supp}(\mathcal F)\subset Z$. En efecto, aquella implica esta
última, puesto que el núcleo de $\rho_{X/Z}$ es entonces igual a $\mathcal F$.
Recíprocamente, si $\operatorname{Supp}(\mathcal F)\subset Z$, basta aplicar
la segunda parte de (5.9.11) al único homomorfismo
$u:\mathcal F\to0$ para concluir que el homomorfismo correspondiente
$v:\mathcal H^0_{X/Z}(\mathcal F)\to0$ es un isomorfismo.
\item[(iii)] Los desarrollos precedentes siguen teniendo sentido para todo
espacio anillado localmente noetheriano en el que toda parte cerrada
irreducible admite exactamente un punto genérico. En particular, se aplican,
sobre un preesquema localmente noetheriano, a cualesquiera \emph{haces de
grupos abelianos} (considerados como módulos sobre el haz simple asociado al
prehaz constante $\mathbf Z$). Para todo $x\in X$ se tiene asimismo el
isomorfismo canónico (5.9.6.1), en el que
$\widetilde{\mathcal F}_x$ designa el haz \emph{inducido} sobre el subespacio
$X_x$ de $X$ por el haz $\mathcal F$; la demostración directa resulta
inmediatamente de la definición (5.9.1.2) y del teorema del doble límite
inductivo.
\end{enumerate}
\end{remarks}

\subsection{\texorpdfstring{Propiedad $(S_2)$ y $Z$-clausura.}{Propiedad (S\string_2) y Z-clausura}}

